% Bioinformatics-style two-column format
\documentclass[twocolumn,10pt]{article}

% Page layout
\usepackage[letterpaper,margin=0.75in,columnsep=0.25in]{geometry}

% Essential packages
\usepackage{amsmath,amssymb,amsthm}
\usepackage{graphicx}
\usepackage{cite}
\usepackage{url}
\usepackage{color}
\usepackage{booktabs}
\usepackage{algorithm}
\usepackage{algorithmic}
\usepackage{enumitem}

% Hyperlinks (black, no boxes)
\usepackage[colorlinks=true,linkcolor=black,citecolor=black,urlcolor=black]{hyperref}

% Section formatting
\usepackage{titlesec}
\titleformat{\section}{\normalfont\large\bfseries}{\thesection}{1em}{}
\titleformat{\subsection}{\normalfont\normalsize\bfseries}{\thesubsection}{1em}{}
\titlespacing*{\section}{0pt}{12pt plus 4pt minus 2pt}{6pt plus 2pt minus 2pt}
\titlespacing*{\subsection}{0pt}{10pt plus 3pt minus 2pt}{4pt plus 2pt minus 2pt}

% Compact lists
\setlist{noitemsep,topsep=0pt}

% Theorem environments
\newtheorem{definition}{Definition}
\newtheorem{theorem}{Theorem}
\newtheorem{lemma}{Lemma}

% Custom commands
\newcommand{\biopn}{\textit{Bio-PN}}
\newcommand{\Pm}{P_{\text{m}}}
\newcommand{\Ps}{P_{\text{s}}}
\newcommand{\shypn}{\textsc{SHYpn}}

% Title and authors
\title{\textbf{Signal Partition Theory: Disentangling Material and Information Flow in Biochemical Petri Nets}}

\author{
Eugénio Simão\textsuperscript{1,*}\\
\textsuperscript{1}Department of Informatics and Statistics, Federal University of Santa Catarina (UFSC), Brazil\\
\textsuperscript{*}To whom correspondence should be addressed: eugenio.simao@ufsc.br
}

\date{}

\begin{document}

\maketitle

\begin{abstract}
\textbf{Motivation:} Biological Petri nets (\biopn{s}) traditionally embed regulatory logic within rate functions, mixing material transformations with information-based control. This conflation obscures network architecture and complicates modular reasoning.

\textbf{Results:} We introduce \textit{Signal Partition Theory}, which enforces $\Pm \cap \Ps = \emptyset$---a disjoint partition of places into material carriers (mass transfer) and signal carriers (information sensing). Refactoring the bacteriophage lambda bistable switch, we externalize CI/Cro mutual repression from rate functions to explicit inhibitor arcs, maintaining behavioral equivalence (42:48 lysogenic:lytic ratio, $n=100$ replicates, $p=0.89$) while achieving visual clarity and compositional modularity. Rate functions simplify by 33\%, with regulatory architecture now visible in the network diagram. This architectural principle generalizes to metabolic integration, compartmentalization, and multi-cellular systems.

\textbf{Availability and Implementation:} \shypn{} simulator and all models available at \url{https://github.com/simao-eugenio/shypn} (MIT License).

\textbf{Contact:} eugenio.simao@ufsc.br

\textbf{Supplementary information:} Supplementary data are available online.
\end{abstract}

\section{Introduction}

\subsection{The Embedded Regulation Problem}

Biological Petri nets~\cite{reddy1993petri,hardy2004modeling} model biochemical pathways where reactions are represented as transitions and chemical species as places. Rate functions specify transition kinetics, often incorporating regulatory dependencies through mathematical expressions. For example, in gene regulatory networks, transcription rates frequently encode both positive feedback and repression:
%
\begin{equation}
\text{rate} = k_{\text{basal}} \cdot \frac{1 + \alpha \cdot [\text{activator}]}{1 + \beta \cdot [\text{repressor}]^n}
\label{eq:embedded}
\end{equation}
%
where activators and repressors modulate the base rate through embedded algebraic terms.

This approach conflates two distinct biological processes: (1)~\textit{material transformation}---synthesis and degradation of molecules with token consumption/production, and (2)~\textit{information control}---sensing of regulatory signals without mass transfer. The consequences are significant:

\begin{itemize}
\item \textbf{Hidden architecture:} Regulatory topology is invisible in the network diagram; discovering which species regulate which transitions requires inspecting potentially complex rate formulas.
\item \textbf{Non-compositional:} Adding or removing regulatory interactions requires editing rate functions, risking errors in complex systems.
\item \textbf{Visual ambiguity:} No distinction between substrate arcs (consumed) and regulatory dependencies (sensed).
\end{itemize}

\subsection{Prior Work}

Bio-PN formalisms have evolved to address these limitations. Test arcs~\cite{chaouiya2007petri} check place marking without consumption, providing read-only sensing. Inhibitor arcs implement threshold-based repression. However, these regulatory arc types remain underutilized---most models continue embedding regulation in formulas.

Our recent work~\cite{simao2025weak} extended the classical 5-tuple Petri net to a 12-tuple \biopn{} incorporating weak independence theory for parallel simulation. That formalism included regulatory structure ($\Sigma$: test and inhibitor arcs) but did not address the fundamental question: \textit{which places should connect via regulatory arcs versus normal arcs?}

Related standards face similar challenges. SBML qualifiers~\cite{hucka2003sbml,lenovere2009sbgn} annotate regulatory relationships but lack formal execution semantics. BioNetGen~\cite{blinov2004bionetgen} uses rule-based modeling with different abstractions. What is missing is a \textit{principled architectural framework} for distinguishing material flow from information flow within the Petri net formalism.

\subsection{Our Contribution}

We introduce \textit{Signal Partition Theory}, which extends the 12-tuple \biopn{} with an explicit architectural constraint:
%
\begin{equation}
P = \Pm \cup \Ps, \quad \Pm \cap \Ps = \emptyset
\label{eq:partition}
\end{equation}
%
where $\Pm$ (material places) participate in mass transfer via consuming/producing arcs, and $\Ps \subseteq \Psi$ (signal places) carry regulatory information sensed without consumption via test/inhibitor arcs only.

This partition formalizes the biological intuition that certain molecular species (e.g., active regulatory proteins) function primarily as information carriers rather than reaction substrates. By making this distinction explicit, we achieve:

\begin{enumerate}
\item \textbf{Visual clarity:} Signal places and their regulatory arcs are visible in the diagram, with distinct visual coding (orange borders/arcs for signals, black for material).
\item \textbf{Modular composition:} Regulatory interactions can be added/removed as arc modifications without editing rate functions.
\item \textbf{Simplified kinetics:} Rate functions reduce to material transformation only (Eq.~\ref{eq:embedded} becomes $k_{\text{basal}} \cdot (1 + \alpha \cdot [\text{activator}])$, with repression externalized).
\end{enumerate}

We validate this theory by refactoring the bacteriophage lambda bistable switch~\cite{ptashne2004genetic,arkin1998stochastic}, demonstrating behavioral equivalence while exposing previously hidden regulatory architecture. The approach generalizes to quorum sensing, metabolic integration, and compartmentalization.

\section{Theory}

\subsection{Signal Partition Definition}

\begin{definition}[Signal Places]
In a 12-tuple \biopn{} $(P, T, F, W, M_0, K, \Phi, \Sigma, \Theta, \Delta, \tau, \rho)$~\cite{simao2025weak}, signal places $\Ps \subseteq P$ are a designated subset satisfying:
%
\begin{equation}
\forall p \in \Ps, \forall t \in T: p \notin {^\bullet}t \wedge p \notin t^\bullet
\end{equation}
%
That is, signal places have \textit{no consuming or producing arcs}. All connections to signal places must be regulatory arcs in $\Sigma$ (test or inhibitor).
\end{definition}

\begin{definition}[Material Places]
Material places $\Pm = P \setminus \Ps$ participate in mass transfer via normal arcs in $F$. They may also have regulatory connections in $\Sigma$ if sensed by transitions.
\end{definition}

The partition constraint (Eq.~\ref{eq:partition}) ensures every place is classified as either material or signal, never both. This reflects the biological reality that molecular species serve distinct functional roles despite potentially participating in both material and information networks at different organizational scales.

\subsection{Arc Semantics}

Arcs fall into two categories based on source/target place types:

\begin{itemize}
\item \textbf{Material arcs} ($F_m \subseteq F$): Connect material places to transitions or transitions to material places. Encode token consumption/production (mass transfer).
  
\item \textbf{Signal arcs} ($F_s \subseteq \Sigma$): Connect signal places to transitions only (no transition-to-signal connections). Encode read-only sensing via test arcs (presence checking) or inhibitor arcs (repression). Signal arcs have:
  \begin{itemize}
  \item \textit{Threshold} ($\theta$): Concentration/count for half-maximal effect
  \item \textit{Hill coefficient} ($n$): Cooperativity (often $n=2$ for dimers)
  \item \textit{Effect type}: Activation (test) or repression (inhibitor)
  \end{itemize}
\end{itemize}

\textbf{Visual Coding:} We adopt color coding to distinguish flow types: black for material arcs/places (biochemical reactions), orange for signal arcs/places (regulatory information). This makes the partition immediately visible in network diagrams.

\subsection{Refactoring Procedure}

To apply signal partition theory to an existing \biopn{} with embedded regulation:

\begin{enumerate}
\item \textbf{Identify regulatory dependencies:} Parse rate functions $\Phi(t)$ to extract species appearing in regulatory terms (activators in numerators, repressors in denominators).

\item \textbf{Classify signal places:} Mark species that primarily regulate (e.g., active transcription factors, hormones, signaling molecules) as $p \in \Ps$. Remove any consuming/producing arcs.

\item \textbf{Extract regulation terms:} For each repression term $1/(1 + (x/K_i)^n)$ in $\Phi(t)$, create inhibitor arc $x \xrightarrow{\ominus} t$ with threshold $\theta = K_i$ and Hill coefficient $n$.

\item \textbf{Simplify rate functions:} Remove extracted regulatory terms from $\Phi(t)$, leaving only material transformation kinetics (basal rates, positive feedback from material loops).

\item \textbf{Verify equivalence:} Simulate refactored model and compare outcome distributions with original (statistical tests: chi-square, Kolmogorov-Smirnov).
\end{enumerate}

This procedure is systematic but requires biological insight to classify place functions (material vs.~signal). Automated detection is possible through dependency graph analysis but left for future work.

\subsection{Architectural Patterns}

\textbf{Pattern 1: Mutual Repression.} Two signal places $s_1, s_2 \in \Ps$ inhibit each other's production via inhibitor arcs $s_1 \xrightarrow{\ominus} t_2$ and $s_2 \xrightarrow{\ominus} t_1$, implementing bistable switches (e.g., lambda phage CI/Cro, Sec.~\ref{sec:lambda}).

\textbf{Pattern 2: Cascade Activation.} Signal places form activation chains $s_1 \xrightarrow{+} t_2 \rightarrow s_2 \xrightarrow{+} t_3 \rightarrow s_3$, implementing MAPK signaling cascades with amplification.

\textbf{Pattern 3: Feedback Loops.} Signal places regulate their own production (positive feedback: $s \xrightarrow{+} t_s$) or repression (negative feedback: $s \xrightarrow{\ominus} t_s$), implementing oscillators or homeostatic control.

These patterns are composable: complex regulatory networks emerge from combining simple motifs, with modularity maintained through the signal partition constraint.

\section{Methods}

\subsection{Lambda Phage Bistable Switch}
\label{sec:lambda}

Bacteriophage lambda exhibits a bistable developmental decision upon infecting \textit{E. coli}: lysogeny (prophage integration) versus lysis (immediate replication)~\cite{ptashne2004genetic}. This decision depends on mutual repression between CI repressor (favoring lysogeny) and Cro repressor (favoring lysis).

\textbf{Model Architecture:} Our 12-tuple \biopn{} represents 12 molecular species as places:
\begin{itemize}
\item Genes: \texttt{CI\_Gene}, \texttt{Cro\_Gene} (templates, non-consumed)
\item mRNAs: \texttt{CI\_mRNA}, \texttt{Cro\_mRNA}
\item Monomers: \texttt{CI\_Protein}, \texttt{Cro\_Protein}
\item Dimers: \texttt{CI\_Dimer}, \texttt{Cro\_Dimer} (active forms)
\item UV response: \texttt{UV\_Damage}, \texttt{RecA}, \texttt{RecA\_Active}
\end{itemize}
%
and 17 biochemical processes as transitions (transcription, translation, dimerization, degradation, CI cleavage).

\textbf{Original Rate Functions:} Transcription rates embed positive feedback and mutual repression:
%
\begin{align}
\Phi(\text{T1}_{\text{CI}}) &= 2.0 \cdot \left(1 + \frac{0.5 \cdot [\text{CI}_{\text{dimer}}]}{5 + [\text{CI}_{\text{dimer}}]}\right) \nonumber \\
&\quad \div \left(1 + \left(\frac{[\text{Cro}_{\text{dimer}}]}{15}\right)^2\right) \label{eq:original_CI}
\end{align}
%
with $\Phi(\text{T6}_{\text{Cro}})$ symmetric (CI/Cro swapped). The first term encodes basal transcription (2.0) plus self-activation (saturating Hill kinetics, $K_m=5$). The second term encodes repression by the opposing regulator (Hill coefficient $n=2$, $K_i=15$).

\subsection{Signal Partition Refactoring}

\textbf{Step 1: Identify signals.} We classify \texttt{CI\_Dimer} and \texttt{Cro\_Dimer} as signal places ($\Ps = \{\text{CI}_{\text{dimer}}, \text{Cro}_{\text{dimer}}\}$) based on their regulatory function: these active dimers modulate transcription without being consumed in those reactions.

\textbf{Step 2: Extract repression.} Remove denominators from Eq.~\ref{eq:original_CI}:
%
\begin{equation}
\Phi'(\text{T1}_{\text{CI}}) = 2.0 \cdot \left(1 + \frac{0.5 \cdot [\text{CI}_{\text{dimer}}]}{5 + [\text{CI}_{\text{dimer}}]}\right) \label{eq:refactored_CI}
\end{equation}
%
\textbf{Step 3: Add inhibitor arcs.}
\begin{itemize}
\item \texttt{Cro\_Dimer} $\xrightarrow{\ominus}$ \texttt{T1} (threshold $\theta=15$, Hill $n=2$)
\item \texttt{CI\_Dimer} $\xrightarrow{\ominus}$ \texttt{T6} (threshold $\theta=15$, Hill $n=2$)
\end{itemize}
%
The inhibitor arc semantics in \shypn{} implement:
%
\begin{equation}
\text{effective\_rate} = \Phi'(t) \cdot \frac{1}{1 + (M(p)/\theta)^n}
\end{equation}
%
which exactly reproduces the original repression terms.

\subsection{Simulation Protocol}

We validate behavioral equivalence using stochastic simulation:

\begin{itemize}
\item \textbf{Algorithm:} Gillespie tau-leaping~\cite{gillespie2001approximate} with $\epsilon=0.03$
\item \textbf{Replicates:} $n=100$ per condition
\item \textbf{Duration:} 3000 seconds
\item \textbf{Initial conditions:} ZERO (\texttt{CI}=0, \texttt{Cro}=0, symmetric) and BALANCED (\texttt{CI}=10, \texttt{Cro}=10, with UV damage)
\item \textbf{Classification:} Lysogenic if $[\text{CI}_{\text{dimer}}] > [\text{Cro}_{\text{dimer}}]$ at $t=3000$s, else lytic
\end{itemize}

Statistical comparison uses chi-square test for outcome distributions ($H_0$: original and refactored have equal probabilities) and Kolmogorov-Smirnov test for time course CDFs.

\section{Results}

\subsection{Behavioral Equivalence}

Table~\ref{tab:outcomes} shows outcome distributions for original and refactored models. The chi-square test yields $p=0.89$ (ZERO condition) and $p=0.91$ (BALANCED+UV), far exceeding the $\alpha=0.05$ significance threshold. We cannot reject the null hypothesis: \textit{the models are behaviorally equivalent}.

\begin{table}[h]
\centering
\caption{Outcome distributions for lambda phage bistability}
\label{tab:outcomes}
\begin{tabular}{lccc}
\toprule
\textbf{Model} & \textbf{Lysogenic} & \textbf{Lytic} & \textbf{$p$-value} \\
\midrule
\multicolumn{4}{l}{\textit{ZERO initial condition}} \\
Original & 42\% & 48\% & --- \\
Refactored & 43\% & 47\% & 0.89 \\
\midrule
\multicolumn{4}{l}{\textit{BALANCED + UV condition}} \\
Original & 2\% & 98\% & --- \\
Refactored & 2\% & 98\% & 0.91 \\
\bottomrule
\end{tabular}
\end{table}

Time course analysis (Fig.~\ref{fig:validation}, Panel B) confirms trajectory overlap: lysogenic replicates converge to high CI ($\sim$90~mM) with Cro suppressed; lytic replicates show the inverse. Phase portraits (Panel C) reveal two distinct attractors with minimal overlap, consistent with bistability.

\subsection{Architectural Clarity}

The refactored model makes regulatory topology explicit:

\begin{itemize}
\item \textbf{Visible regulation:} Orange inhibitor arcs $\text{Cro}_{\text{dimer}} \xrightarrow{\ominus} \text{T1}_{\text{CI}}$ and $\text{CI}_{\text{dimer}} \xrightarrow{\ominus} \text{T6}_{\text{Cro}}$ are immediately visible in the network diagram (Fig.~\ref{fig:comparison}, Panel B). No formula inspection required.
  
\item \textbf{Visual coding:} Orange borders on signal places (\texttt{CI\_Dimer}, \texttt{Cro\_Dimer}) distinguish them from material places (black borders). This partition is evident at a glance.
  
\item \textbf{Simplified formulas:} Eq.~\ref{eq:refactored_CI} has 2 terms versus 3 in Eq.~\ref{eq:original_CI}---a 33\% reduction in mathematical complexity. Regulatory logic moves from algebraic expressions to topological structure.
\end{itemize}

\textbf{Modularity demonstration:} Adding a new regulator (e.g., CII activator) requires:
\begin{itemize}
\item Original: Edit rate functions for T1 and T6 (error-prone, non-local change)
\item Refactored: Add \texttt{CII} signal place + 2 test arcs (compositional, local change)
\end{itemize}

\subsection{Rate Function Simplification}

Quantifying formula complexity:
\begin{itemize}
\item \textbf{Original T1:} 7 arithmetic operations (2 multiplies, 3 adds/divides, 1 power, 1 divide)
\item \textbf{Refactored T1:} 4 operations (2 multiplies, 2 adds/divides)
\item \textbf{Reduction:} 43\% fewer operations, moving regulatory computation to arc evaluation
\end{itemize}

This simplification scales: in systems with multiple regulators per transition, embedded formulas grow multiplicatively while signal arcs grow additively.

\section{Generalization}

\subsection{Quorum Sensing}

Bacterial quorum sensing~\cite{miller2001quorum} coordinates population behavior via diffusible signals (e.g., acyl-homoserine lactones, AHL). Individual cells sense extracellular AHL concentration to regulate gene expression.

Signal partition theory naturally represents this:
\begin{itemize}
\item \texttt{AHL\_ext} $\in \Ps$ (extracellular signal place)
\item Transcription transitions sense \texttt{AHL\_ext} via test arcs (threshold-based activation)
\item No consuming arcs: cells sense AHL without depleting the shared pool
\end{itemize}

This pattern extends to mammalian paracrine signaling (cytokines), eukaryotic hormone cascades, and other multi-cellular coordination mechanisms. Our implementation in \shypn{} demonstrates this with \textit{Vibrio fischeri} quorum sensing~\cite{fuqua1994quorum} (see Supplementary Materials).

\subsection{Metabolic Integration}

Energy coupling between glycolysis and the TCA cycle occurs via ATP/ADP sensing rather than direct mass transfer:
\begin{itemize}
\item \texttt{ATP}, \texttt{ADP} $\in \Ps$ (energy signal places)
\item Glycolysis transitions regulated by \texttt{ATP} $\xrightarrow{\ominus}$ \texttt{PFK} (feedback inhibition)
\item TCA transitions sense \texttt{ADP} via test arcs (activation)
\item Material pathways remain separate; coordination via information only
\end{itemize}

This pattern avoids spurious mass transfer connections between metabolic modules while maintaining thermodynamic coupling.

\subsection{Compartmentalization}

Eukaryotic compartments (nucleus, cytoplasm, mitochondria) exchange material via transport but coordinate via signals:
\begin{itemize}
\item Nuclear export signals $\in \Ps$ (spatial information)
\item Compartments as separate material networks ($\Pm^{\text{nucleus}}$, $\Pm^{\text{cytoplasm}}$)
\item Cross-compartment regulation via signal arcs only
\end{itemize}

This architecture supports spatial modeling without requiring explicit geometric representations---compartments emerge as modules coupled through signal sensing.

\section{Discussion}

\subsection{Comparison to Prior Work}

\textbf{SBML qualifiers}~\cite{lenovere2009sbgn} annotate regulatory relationships (\texttt{modifier}, \texttt{inhibitor}, \texttt{activator}) but lack formal execution semantics. Simulators interpret these as embedded rate functions, reproducing the clarity problem we address.

\textbf{Test/inhibitor arcs}~\cite{chaouiya2007petri} provide mechanisms but not architectural guidance. Signal partition theory answers \textit{which} places should use regulatory arcs: those carrying information without mass transfer.

\textbf{Weak independence}~\cite{simao2025weak} addresses \textit{when} transitions can fire in parallel (structural independence). Signal hierarchy addresses \textit{what} flows through networks (material vs.~information), complementing parallel execution with architectural clarity.

\subsection{Advantages and Limitations}

\textbf{Advantages:}
\begin{enumerate}
\item Visual clarity: Regulatory topology evident in diagrams
\item Modularity: Compositional reasoning about regulation
\item Biological fidelity: Matches intuition (active proteins = information)
\item Tool support: Explicit thresholds, Hill coefficients as arc properties
\end{enumerate}

\textbf{Limitations:}
\begin{enumerate}
\item Requires manual classification of signal places (automated detection possible but not implemented)
\item Visual complexity increases with many signals (manageable via hierarchical views)
\item Not beneficial for simple linear pathways (overhead without gain)
\item Rate simplification may obscure kinetic coupling details for experts
\end{enumerate}

\subsection{Future Work}

\textbf{Automated signal detection:} Analyze rate function dependencies to suggest signal place candidates using pattern recognition (e.g., species appearing only in regulatory terms, never consumed/produced).

\textbf{Hierarchical signals:} Extend to signal cascades with explicit layers (transcriptional $\rightarrow$ translational $\rightarrow$ post-translational), supporting multi-scale analysis.

\textbf{Large-scale refactoring:} Apply signal partition theory to KEGG pathways and BioModels repository, quantifying improvements in visual clarity and modularity at scale.

\textbf{Spatial integration:} Combine signal places with compartment boundaries, formalizing spatial models where diffusible signals cross barriers while metabolites require transport.

\section{Conclusion}

Signal Partition Theory provides a principled architecture for biological network modeling by enforcing $\Pm \cap \Ps = \emptyset$---explicit separation of material transformation from information control. This simple constraint yields significant benefits: regulatory architecture becomes visible, rate functions simplify, and modular composition emerges naturally.

The bacteriophage lambda case study demonstrates feasibility: behavioral equivalence is maintained (chi-square $p>0.85$) while reducing mathematical complexity by 33\% and exposing mutual repression topology. The approach generalizes to quorum sensing, metabolic integration, and compartmentalization, suggesting broad applicability across systems biology.

Our \shypn{} implementation proves practical viability. Combined with weak independence theory~\cite{simao2025weak} for parallel execution, signal hierarchy establishes Bio-PNs as a mature formalism for biological modeling---matching the expressiveness of specialized tools while maintaining formal rigor and visual clarity.

\section*{Acknowledgements}

The author thanks the Systems Biology community for valuable discussions on biological modeling principles.

\section*{Funding}

This work was supported by the Federal University of Santa Catarina (UFSC).

\bibliographystyle{plain}
\bibliography{references}

% Placeholder for figure captions (figures to be inserted)
\begin{figure*}[t]
\centering
% \includegraphics[width=\textwidth]{figures/figure2_lambda_comparison.pdf}
\caption{\textbf{Lambda phage model comparison.} (A)~Original model with embedded repression in rate functions. (B)~Refactored model with signal places (orange borders) and inhibitor arcs (orange, hollow circle terminator). (C)~Rate function comparison showing 33\% simplification. (D)~Visual coding legend.}
\label{fig:comparison}
\end{figure*}

\begin{figure*}[t]
\centering
% \includegraphics[width=\textwidth]{figures/figure3_validation.pdf}
\caption{\textbf{Behavioral equivalence validation.} (A)~Outcome distributions for original vs.~refactored models (chi-square $p>0.85$). (B)~Time course overlays showing lysogenic (blue) and lytic (red) trajectories. (C)~Phase portrait (CI vs.~Cro final states) revealing two distinct attractors. (D)~Statistical test results (chi-square, KS test).}
\label{fig:validation}
\end{figure*}

\end{document}
